Everything that follows rests on treating patch tokenisation as an operator rather than an implementation detail, because its degeneracies are then derivable instead of observed. This section fixes that operator and the notation the rest of the report uses for it; the conditions under which it loses information are derived in \autoref{sec:two}.

Let $\mathbf{x} = \{x[n]\}_{n \in \mathbb{N}}$, with $\mathbb{N}=\{0,1,\ldots\}$, be a discrete-time signal representing the input time series. Defining the patch-based tokenisation operator $\mathcal{T}_{P,S}$, with patch size $P \in \mathbb{N}^{+}$ and stride $S \in \mathbb{N}^{+}$, as a mapping that transforms the signal $\mathbf{x}$ into a sequence of vector-valued tokens $\mathbf{T}=\{\mathbf{t}_k\}_{k\in\mathbb{N}}$ with $\mathbf{t}_k\in\mathbb{R}^{P}$, the $k$-th token is

\begin{equation*}
\mathbf{t}_k
=
\mathcal{T}_{P,S}(\mathbf{x})[k]
=
\begin{bmatrix}
x[kS] \\
x[kS+1] \\
\vdots \\
x[kS+P-1]
\end{bmatrix}.
\end{equation*}

Using element-wise indexing notation,

\begin{equation}
\label{eq:element_def}
\mathbf{t}_k[m]
=
x[kS+m],
\qquad
m\in\{0,\ldots,P-1\}.
\end{equation}
